\section{Jacobson Radical of a Ring}

\begin{definition}[Jacobson radical]
    Let $R$ be a ring with identity. The \emph{Jacobson radical} of $R$ is defined by
    \[
        J(R):=\bigcap \{\, M \mid M \text{ is a maximal left ideal of } R \,\}.
    \]
\end{definition}

\begin{remark}[Annihilators]
    Let ${}_RM$ be a left $R$-module.

    \begin{enumerate}
        \item For $x\in M$, define
              \[
                  \operatorname{Ann}_R(x):=\{\,r\in R \mid rx=0\,\}.
              \]
              This is a left ideal of $R$.

        \item For a subset $X\subseteq M$, define
              \[
                  \operatorname{Ann}_R(X):=\bigcap_{x\in X}\operatorname{Ann}_R(x).
              \]

        \item If $N\leq M$ is a submodule, then
              \[
                  \operatorname{Ann}_R(N):=\bigcap_{x\in N}\operatorname{Ann}_R(x)
              \]
              is a two-sided ideal of $R$.
    \end{enumerate}
\end{remark}

\begin{theorem}[Properties of the Jacobson radical]
    Let $R$ be a ring with identity. Then the following hold.

    \begin{enumerate}
        \item
              \[
                  J(R)=\bigcap \{\,\operatorname{Ann}_R(S)\mid S \text{ is a simple left }R\text{-module}\,\}.
              \]
              In particular, $J(R)$ is a two-sided ideal of $R$.

        \item For $a\in R$,
              \[
                  a\in J(R)
                  \quad\Longleftrightarrow\quad
                  1-xa \text{ has a left inverse for every } x\in R.
              \]
              Equivalently, for every $x\in R$ there exists $b\in R$ such that
              \[
                  b(1-xa)=1.
              \]

        \item $J(R)$ is the largest two-sided ideal $I$ of $R$ satisfying the condition
              \[
                  a\in I \;\Longrightarrow\; 1-a \text{ is a unit in } R.
              \]

        \item $J(R)$ coincides with the intersection of all maximal right ideals of $R$.
    \end{enumerate}
\end{theorem}

\begin{proof}
    \textbf{(1)} Let $S$ be a simple left $R$-module, and choose $0\neq v\in S$. Since $S$ is simple,
    \[
        S=Rv.
    \]
    Consider the surjective $R$-homomorphism
    \[
        R\longrightarrow S,\qquad r\longmapsto rv.
    \]
    Its kernel is $\operatorname{Ann}_R(v)$, so
    \[
        S\cong R/\operatorname{Ann}_R(v).
    \]
    Because $S$ is simple, $\operatorname{Ann}_R(v)$ is a maximal left ideal of $R$. Hence
    \[
        J(R)\subseteq \operatorname{Ann}_R(v)
    \]
    for every nonzero $v\in S$, and therefore
    \[
        J(R)\subseteq \operatorname{Ann}_R(S)
    \]
    for every simple left $R$-module $S$.

    Conversely, suppose that
    \[
        x\in \bigcap \{\,\operatorname{Ann}_R(S)\mid S \text{ simple left }R\text{-module}\,\}.
    \]
    Let $M$ be any maximal left ideal of $R$. Then $R/M$ is a simple left $R$-module, so
    \[
        x\in \operatorname{Ann}_R(R/M).
    \]
    In particular,
    \[
        x(1+M)=0,
    \]
    hence
    \[
        x+M=0,
    \]
    so $x\in M$. Since this holds for every maximal left ideal $M$, we obtain
    \[
        x\in J(R).
    \]
    Thus
    \[
        J(R)=\bigcap \{\,\operatorname{Ann}_R(S)\mid S \text{ simple left }R\text{-module}\,\}.
    \]

    Finally, each $\operatorname{Ann}_R(S)$ is a two-sided ideal, so their intersection $J(R)$ is also a two-sided ideal.

    \medskip

    \textbf{(2)} Assume first that $a\in J(R)$. Suppose, for contradiction, that there exists $x\in R$ such that $1-xa$ has no left inverse. Then the principal left ideal
    \[
        R(1-xa)
    \]
    is a proper left ideal of $R$, so it is contained in some maximal left ideal $M$.

    Since $a\in J(R)$, we have $a\in M$, and because $M$ is a left ideal,
    \[
        xa\in M.
    \]
    Also,
    \[
        1-xa\in M.
    \]
    Therefore
    \[
        1=(1-xa)+xa\in M,
    \]
    a contradiction. Hence $1-xa$ has a left inverse for every $x\in R$.

    Conversely, assume that for every $x\in R$, the element $1-xa$ has a left inverse. Suppose that $a\notin J(R)$. Then there exists a maximal left ideal $M$ such that
    \[
        a\notin M.
    \]
    Hence
    \[
        M+Ra=R,
    \]
    so there exist $b\in M$ and $x\in R$ such that
    \[
        b+xa=1.
    \]
    Thus
    \[
        1-xa=b\in M.
    \]
    But an element of a proper left ideal cannot have a left inverse, for if $c(1-xa)=1$, then
    \[
        1\in M,
    \]
    contradicting the propriety of $M$. This contradiction shows that $a\in J(R)$.

    \medskip

    \textbf{(3)} First let $a\in J(R)$. By part \textbf{(2)} with $x=1$, there exists $b\in R$ such that
    \[
        b(1-a)=1.
    \]
    Then
    \[
        b=1+ba=1-(-b)a.
    \]
    Since $a\in J(R)$, part \textbf{(2)} again implies that $b$ has a left inverse. So there exists $c\in R$ such that
    \[
        cb=1.
    \]
    Now
    \[
        c=cb(1-a)=1-a.
    \]
    Hence
    \[
        (1-a)b=1,
    \]
    and together with $b(1-a)=1$, this shows that $1-a$ is a unit in $R$.

    Thus $J(R)$ satisfies the stated property.

    Now let $I$ be a two-sided ideal of $R$ such that
    \[
        a\in I \;\Longrightarrow\; 1-a \text{ is a unit in }R.
    \]
    Suppose $I\not\subseteq J(R)$. Then there exists a maximal left ideal $M$ such that
    \[
        I\not\subseteq M.
    \]
    Choose $a\in I\setminus M$. Then
    \[
        M+Ra=R,
    \]
    so there exist $m\in M$ and $r\in R$ such that
    \[
        m+ra=1.
    \]
    Since $I$ is a two-sided ideal and $a\in I$, we have
    \[
        ra\in I.
    \]
    Therefore
    \[
        1-ra
    \]
    is a unit in $R$. But
    \[
        1-ra=m\in M,
    \]
    and a proper left ideal cannot contain a unit. This is a contradiction. Hence
    \[
        I\subseteq J(R).
    \]
    So $J(R)$ is the largest such two-sided ideal.

    \medskip

    \textbf{(4)} Let
    \[
        J_r(R):=\bigcap \{\,N\mid N \text{ is a maximal right ideal of }R\,\}.
    \]
    Applying the preceding arguments to the opposite ring $R^{\operatorname{op}}$, we see that $J_r(R)$ is also the largest two-sided ideal $I$ of $R$ such that
    \[
        a\in I \;\Longrightarrow\; 1-a \text{ is a unit in }R.
    \]
    By part \textbf{(3)}, this property characterizes $J(R)$ uniquely. Therefore
    \[
        J_r(R)=J(R).
    \]
    Hence $J(R)$ coincides with the intersection of all maximal right ideals of $R$.
\end{proof}

\begin{definition}
    Let $R$ be a ring with $1$.

    \begin{enumerate}
        \item An element $a\in R$ is called \emph{nilpotent} if
              \[
                  a^n=0
              \]
              for some integer $n\ge 1$.

        \item A left ideal $I\subseteq R$ is called \emph{nilpotent} if
              \[
                  I^n=0
              \]
              for some integer $n\ge 1$.

        \item A left ideal $I\subseteq R$ is called a \emph{nil left ideal} if every element of $I$
              is nilpotent.
    \end{enumerate}
\end{definition}

\begin{lemma}
    If $a\in R$ is nilpotent, say $a^n=0$, then $1-a$ is a unit in $R$, with inverse
    \[
        (1-a)^{-1}=1+a+a^2+\cdots+a^{n-1}.
    \]
\end{lemma}

\begin{proof}
    Since $a^n=0$, we have
    \[
        (1-a)(1+a+a^2+\cdots+a^{n-1})=1-a^n=1,
    \]
    and similarly
    \[
        (1+a+a^2+\cdots+a^{n-1})(1-a)=1-a^n=1.
    \]
    Hence $1-a$ is a unit.
\end{proof}

\begin{lemma}
    Every nilpotent left ideal is a nil left ideal.
\end{lemma}

\begin{proof}
    Suppose $I^n=0$ and let $a\in I$. Then
    \[
        a^n\in I^n=0,
    \]
    so $a$ is nilpotent. Hence every element of $I$ is nilpotent.
\end{proof}

\begin{remark}
    We shall use the following standard characterization of the Jacobson radical:
    \[
        a\in J(R)
        \quad\Longleftrightarrow\quad
        \text{for every }x\in R,\; 1-xa \text{ has a left inverse in }R.
    \]
    In particular, if $a\in J(R)$, then $1-a$ is a unit.
\end{remark}

\begin{theorem}
    Let $I$ be a nil left ideal of a ring $R$. Then
    \[
        I\subseteq J(R).
    \]
\end{theorem}

\begin{proof}
    Let $a\in I$. For every $x\in R$, since $I$ is a left ideal we have $xa\in I$. Because $I$
    is nil, the element $xa$ is nilpotent. By the previous lemma, $1-xa$ is a unit in $R$, hence
    in particular it has a left inverse.

    Therefore, by the characterization of $J(R)$,
    \[
        a\in J(R).
    \]
    Since $a\in I$ was arbitrary, it follows that
    \[
        I\subseteq J(R).
    \]
\end{proof}

